\subsection{CX-N14: proved cases of splitting a prime-field algebra from an automorphism}
\label{finding:CX-N14}

\paragraph{Finding.}
Let \(A\) be a commutative unital \(\mathbb F_p\)-algebra
given by an explicit basis and multiplication table, promised to
be isomorphic to \(\mathbb F_p^d\).
Let \(\tau\ne1\) be a supplied \(\mathbb F_p\)-algebra automorphism,
given by its matrix, and let \(h=\operatorname{ord}(\tau)\).
If at least one of the following conditions holds, where
\(r\) in (ii) and (iii) is a prime divisor of \(h\),
\[
\begin{array}{ll}
\textnormal{(i)} & p\leq d,\\
\textnormal{(ii)} & r\mid(p-1),\\
\textnormal{(iii)} & r\ne p,\quad
 v_r\!\left(p^{s}-1\right)=1,\quad s=\operatorname{ord}_r(p),
\end{array}
\]
then a nonzero zero divisor of \(A\) can be found deterministically
in \(\operatorname{poly}(d,\log p)\) bit operations.
Cases (i) and (iii) are unconditional.
Case (ii) uses only \(A_D\), through CX-N02.
For \(A=\mathbb F_p[X]/(f)\), the output yields a proper factor of
\(f\).

This is a consequence of the stated constructions and the precise
algebraic algorithms of
\cite[Lemma~4.2 and Proposition~4.6]{CXNIKRS08};
it is not presented as a new discovery of those source algorithms.
The finding identifies concrete parts of the general
automorphism-splitting step whose analytic requirements are
discharged or absent.

\paragraph{Order recovery and reduction to prime order.}
An automorphism of \(\mathbb F_p^d\) permutes the primitive
idempotents, so its order divides
\(T=\operatorname{lcm}(1,\ldots,d)\).
The bit length of \(T\) is polynomial in \(d\), and its prime factors
are all at most \(d\).
Starting with \(T\), divide out a prime factor whenever the
corresponding power of \(\tau\) is the identity.
Matrix powering performs these tests and determines \(h\) in
polynomial time.
The same small-prime factorization finds the prime divisors \(r\)
of \(h\).
Replace \(\tau\) by \(\tau^{h/r}\) for a chosen \(r\).
This has exact order \(r\), and \(r\leq d\).
The integer \(s=\operatorname{ord}_r(p)\), when \(r\ne p\), and
the valuation test in (iii) are also computable in polynomial
time: \(s\leq r-1\), and \(p^s\) has \(O(d\log p)\) bits.
Thus selecting a covered case requires no unprovided integer
factorization oracle.

\paragraph{Proof of case (i).}
Here \(p\leq d\); the supplied automorphism is not needed for
this case. In particular, it covers every characteristic-order
case \(r=p\), since a prime divisor of \(h\) is at most \(d\).
Choose a basis element \(b\) outside the one-dimensional scalar
subspace \(\mathbb F_p1\).
Under the promised product decomposition its coordinates are not
all equal.
For some \(a\in\mathbb F_p\), the element \(b-a1\) has a zero
coordinate, and it is nonzero for every \(a\).
Enumerate the \(p\leq d\) scalars and test the multiplication
matrices of \(b-a1\) for singularity.
A singular matrix gives a nonzero zero divisor.
This requires polynomial time and no analytic assumption.
\hfill\(\square\)

\paragraph{Proof of case (ii).}
Use CX-N02 to find an \(r\)-th nonresidue \(a\in\mathbb F_p^*\).
Set \(\zeta=a^{(p-1)/r}\), of exact order \(r\).
The prime-order permutation \(\tau\) has an \(r\)-cycle, so its
\(\zeta\)-eigenspace is nonzero.
Linear algebra finds
\[
 0\ne z\in A,\qquad \tau(z)=\zeta z.
\]
If multiplication by \(z\) is singular, it is already the desired
zero divisor. Otherwise \(z\) is a unit.

Compute \(B=A^\tau=\ker(\tau-1)\) and its multiplication table by
linear algebra.
It is a split \(\mathbb F_p\)-algebra, with components indexed by
the orbits of \(\tau\).
The unit \(c=z^r\) belongs to \(B\).
On each orbit its constant coordinate is the \(r\)-th power of
some coordinate of \(z\), so \(c\) has an \(r\)-th root in \(B\).

The elementary-algebra root algorithm in
\cite[\S2.4, Lemma~2.6, author pp.7--9]{CXNGao01}
finds such a \(y\in B\) in
\(\operatorname{poly}(d,r,\log p)\) time given a primitive
\(r^e\)-th root of unity, where \(p-1=r^e w\), \(r\nmid w\).
The same nonresidue \(a\) supplies that root as \(a^w\).
Repeated division by \(r\) computes \(e,w\).
This uses no extension-field nonresidue and no factorization of
the complementary factor \(w\).

Now
\[
 \prod_{j=0}^{r-1}(z-\zeta^j y)=z^r-y^r=0.
\]
Every factor is nonzero: \(z=\zeta^j y\) would imply, on applying
\(\tau\), that \(\zeta z=z\), contradicting invertibility of \(z\).
At least one factor is a nonunit, since a product of units cannot
be zero.
Testing the \(r\leq d\) multiplication matrices finds a nonzero
nonunit, hence a nonzero zero divisor.
All operations have polynomial bit complexity.
\hfill\(\square\)

\paragraph{Constructive removal of fixed components.}
For case (iii), and also as an unconditional preliminary test in
any prime-order case, form
\[
 J=\bigl(a-\tau(a):a\in A\bigr)\subseteq A.
\]
For an explicit basis \(b_1,\ldots,b_d\), this ideal is the linear
span of the \(d^2\) products
\(b_j(b_i-\tau(b_i))\), so it is computable by linear algebra.
It is nonzero because \(\tau\ne1\).
On a component fixed by \(\tau\), every generator vanishes.
On a moved component, an abstract primitive idempotent supported
there proves that some generator has nonzero coordinate.
Consequently \(J\) is exactly the ideal supported on the moved
components.

If \(J\ne A\), any nonzero element of \(J\) is a zero divisor,
annihilated by the nonzero complementary component ideal.
This immediately solves the problem.
If \(J=A\), every cycle of \(\tau\) has length \(r\).
The cyclic action is then semiregular in the sense of IKRS:
no nonidentity power acts as the identity on a nonzero ideal.
This establishes the required promise for the next step
constructively; no component decomposition has been assumed as
algorithmic input.

\paragraph{Proof of case (iii).}
We may assume the preceding test gave \(J=A\).
Put
\[
 D_r=\mathbb F_p[T]/(\Phi_r(T)),\qquad
 C=A\otimes_{\mathbb F_p}D_r,\qquad
 \zeta=T\bmod\Phi_r(T).
\]
The algebra \(D_r\) is semisimple because \(p\ne r\).
Every primitive \(r\)-th root has Frobenius orbit length
\(s=\operatorname{ord}_r(p)\).
It follows that \(D_r\) is a product of copies of
\(\mathbb F_{p^s}\), and so is \(C\), since \(A\) is split over
\(\mathbb F_p\).
No explicit factorization of \(\Phi_r\) is needed for this argument
or for the construction of \(C\).

The valuation-one hypothesis implies that every element of
\(r\)-power order in \(C^\times\) has order dividing \(r\).
For \(1\leq b<r\), let \(\rho_b\) fix \(A\) and send
\(\zeta\) to \(\zeta^b\).
In this exponent-\(r\) case, the Teichm\"uller subgroup used in IKRS
is exactly
\[
 \mathcal T_{A,r}
 =\{u\in C^\times:u^r=1,\ 
          \rho_b(u)=u^b\text{ for all }1\leq b<r\};
\]
their Teichm\"uller exponent lift is congruent to \(b\) modulo \(r\).

IKRS Proposition~4.6, with the semiregular action just verified,
constructs \(x\in\mathcal T_{A,r}\) satisfying
\(\tau(x)=\zeta x\) in polynomial time.
The element \(\zeta\) also belongs to \(\mathcal T_{A,r}\).
Both have order \(r\).
Moreover \(x\notin\langle\zeta\rangle\), since \(\tau\) fixes
\(\zeta\) but does not fix \(x\).
Thus \(\langle x,\zeta\rangle\cong C_r\times C_r\) is noncyclic.
IKRS Lemma~4.2 now produces a nonzero zero divisor in
\(A\), in \(\operatorname{poly}(r,\log|A|)\) time.
Its conclusion is in \(A\) itself: the source proof makes an
annihilator idempotent invariant under every \(\rho_b\), and
\(C^{\{\rho_b\}}=A\).
Since \(r\leq d\) and \(\log|A|=d\log p\), the required bit
complexity follows, without GRH.
\hfill\(\square\)

\paragraph{The remaining branch of this argument.}
Under \(A_D\), if none of the three cases applies, every relevant
prime-order branch has
\[
 p>d,\qquad
 r\nmid p(p-1),\qquad
 s=\operatorname{ord}_r(p)>1,\qquad
 r^2\mid p^s-1.
\]
The fixed-component test may still split individual instances.
For a semiregular instance in this branch, the above proof supplies
the element \(x\), but does not supply a coefficient
Teichm\"uller unit of sufficient \(r\)-power order in \(D_r\),
or an extension-field radical algorithm that replaces it.
Gao's general automorphism-splitting Lemma~3.6 cites
R\'onyai's 1992 result
\cite[Lemma~3.6, author p.16]{CXNGao01};
the exact additional analytic input in a chosen general route
must still be tracked.
The algebra \(A[\zeta_r]\) in IKRS has extension-field components
when \(s>1\).
These statements delimit this proved algorithm.
They do not assert that a Dirichlet-only alternative for the
remaining branch is impossible.

\paragraph{A precise sufficient algebraic input for the remaining branch.}
Let \(e=v_r(p^s-1)\), and suppose an element
\(\eta\in D_r^\times\), in the explicit power basis of \(D_r\),
is additionally supplied with
\[
 \operatorname{ord}(\eta)=r^e,\qquad
 \rho_b(\eta)=\eta^{\omega_e(b)}
 \quad(1\leq b<r),
 \qquad
 \omega_e(b)=b^{r^{e-1}}\bmod r^e.
\]
These promises can be checked by polynomially many algebra
operations; in particular, exact order is checked with the two
powers \(r^e\) and \(r^{e-1}\).
They say that \(\eta\) belongs to the coefficient
Teichm\"uller subgroup, with maximum possible \(r\)-power order.
Embedded in \(C\), it also belongs to \(\mathcal T_{A,r}\)
and is fixed by \(\tau\).
Every \(r\)-element of \(C^\times\), including the element \(x\)
from Proposition~4.6, has order dividing \(r^e\), because each
field component has \(p^s\) elements.
If \(\langle x,\eta\rangle\) were cyclic, its order would be
\(r^e\), and its element \(\eta\) of that order would generate
it. This would put \(x\) in \(\langle\eta\rangle\), contradicting
\(\tau(x)=\zeta x\ne x\).
Thus the group is noncyclic, and Lemma~4.2 gives a nonzero zero
divisor in \(A\) in polynomial time.
This explicitly delimits a sufficient additional input; the
construction of such an \(\eta\) in the remaining branch is not
claimed here under \(A_D\).
\hfill\(\square\)

\paragraph{A limit of the valuation-one argument.}
The valuation hypothesis in CX-N14 cannot simply be deleted from
its particular noncyclicity proof.
For an explicit example, take \(p=17\), \(r=3\), and
\(A=\mathbb F_{17}^3\).
The polynomial \(T^2+T+1\) has no root over \(\mathbb F_{17}\),
since a root would have order three in its multiplicative group
of order sixteen. Thus
\[
 D=\mathbb F_{17}[T]/(T^2+T+1)\cong\mathbb F_{17^2},
 \qquad \zeta=T\bmod(T^2+T+1).
\]
Here \(s=\operatorname{ord}_3(17)=2\) and
\(v_3(17^2-1)=v_3(288)=2\).
The cyclic group \(D^\times\) contains an element \(\eta\) of
order nine with \(\eta^3=\zeta\): start with a generator of its
order-nine subgroup and replace it by its inverse if necessary.
Set
\[
 x=(\eta,\zeta\eta,\zeta^2\eta)\in
 A\otimes_{\mathbb F_{17}}D=D^3,
 \qquad
 \tau(a_0,a_1,a_2)=(a_1,a_2,a_0).
\]
Then \(\tau(x)=\zeta x\), \(x\) has order nine, and
\(x^3=(\zeta,\zeta,\zeta)\).
The nonidentity cyclotomic automorphism \(\rho_2\) of \(D\)
is its seventeenth-power Frobenius.
On every element of order dividing nine, this acts by exponent
\(17\equiv-1\pmod9\), which is also
\(\omega_2(2)=2^3\equiv-1\pmod9\).
It follows coordinate by coordinate that
\(\rho_2(x)=x^{\omega_2(2)}\); hence \(x\) is in the required
Teichm\"uller subgroup.
Nevertheless \(\langle x,\zeta\rangle=\langle x\rangle\)
is cyclic of order nine.
This example satisfies the eigenunit properties used from
Proposition~4.6 but does not force that pair to be noncyclic.
It makes no assertion about which valid eigenunit a particular
implementation returns, or about the impossibility of a different
Dirichlet-only algorithm.
\hfill\(\square\)

\paragraph{Recovering a polynomial factor.}
If \(A=\mathbb F_p[X]/(f)\), represent the resulting nonzero zero
divisor by \(g\) of degree less than \(\deg f\).
Then \(\gcd(f,g)\) has positive degree because multiplication by
\(g\) is singular, and has degree less than \(\deg f\) because
\(g\ne0\) in \(A\). The Euclidean algorithm therefore returns
a proper factor in polynomial time.

\paragraph{Geometric restatement.}
The algebra \(A\) is the ring of functions on a constant finite
\'etale scheme of degree \(d\).
Its automorphism permutes the components.
The ideal \(J\) detects the moved-component open-and-closed
subscheme; the fixed algebra \(A^\tau\) is the orbit quotient.
In case (ii), a prime-field eigenfunction and an invariant radical
yield a proper component decomposition.
In case (iii), the cyclotomic coefficient extension admits the
stated Teichm\"uller construction, and the source descent returns
an idempotent on the original finite \'etale scheme.
The explicit algebra encoding and divisibility/valuation conditions
remain part of this geometric statement.
